\section{Related work}
\label{sec:related-work}

The published work that bears on the aggregation of capital across risk types falls into five groups: the additive and the
square-root rules together with the supervisory view of them, the copula models of the dependence between risk types, the sharp
bounds that hold whatever the dependence, the theory of subadditivity that explains why the bounds are not where intuition puts
them, and the measures of model risk. The proposed protocol uses one or more members of each group, and the last subsection says
what none of them does on its own.

\subsection{Additive and square-root aggregation, and the supervisory view}
\label{sec:rw-additive}

The additive rule sums the stand-alone capital of each risk type and so assumes that the worst outcomes of the types coincide.
The square-root rule, also called the variance-covariance formula, combines the stand-alone figures through a matrix of
correlations; it is exact for elliptical losses and is otherwise a convention. Kuritzkes, Schuermann and Weiner studied risk
measurement and capital adequacy in financial conglomerates, where the question is how much the combination of risk types
reduces the capital of the group \citep{kuritzkes2003}. Alexander and Pezier addressed the aggregation of firm-wide market and
credit risks \citep{alexander2004}. Rosenberg and Schuermann built an integrated model of market, credit and operational risk in
which each marginal is skewed and fat-tailed and the dependence is a copula; they report that the additive approximation
overestimates the risk of a typical large bank by about 30 to 40 percent and that the aggregate is sensitive to the copula and to
the correlations \citep{rosenberg2006}. Dimakos and Aas, and Aas, Dimakos and \O ksendal, described the building of a total
economic capital figure from the loss distributions of the risk types by simulation with a stated dependence between them
\citep{dimakos2004,aas2007}. Brockmann and Kalkbrener treat the aggregation of risk in an economic capital framework
\citep{brockmann2010}.

The supervisory literature is cautious about all of these. The Basel Committee reviewed the range of practices and the issues in
economic capital frameworks \citep{bcbs2009}, and the Joint Forum found that supervisors generally do not rely on the aggregation
models of firms, which it described as a work in progress \citep{jointforum2010}. The common finding of this group is that the
aggregate depends on a choice, the rule or the matrix, that the data seldom settle. None of these works prescribes that the
dependence of the aggregate be displayed as a set of figures next to the one figure that is used.

\subsection{Copula models of the dependence between risk types}
\label{sec:rw-copulas}

A copula separates the marginal distribution of each risk from the dependence between them, which is the structure the aggregation
needs: the marginals come from the models of the risk types and the copula is the one added assumption. Embrechts, McNeil and
Straumann set out the properties and the pitfalls of the linear correlation as a measure of dependence, among them that the
marginals and the correlation matrix do not determine the joint law and that the value at risk of a sum need not be subadditive
\citep{ems2002}; Embrechts, Lindskog and McNeil give the copula framework with its applications in risk management
\citep{elm2003}. Li introduced the Gaussian copula in credit risk \citep{li2000}, and Demarta and McNeil describe the $t$ copula
and its relatives, in particular the tail dependence that the $t$ copula has and the Gaussian copula lacks \citep{demarta2005}.
Frey and McNeil compare models of dependent defaults in portfolio credit risk by the dependence they imply, a comparison that
matters for the credit marginal of the aggregate \citep{frey2003}. Lindskog, McNeil and Schmock show that Kendall's rank
correlation of an elliptical copula is a function of its correlation parameter and does not depend on the degrees of freedom
\citep{lindskog2003}; the consequence, used in Section~\ref{sec:theory}, is that a Gaussian and a $t$ copula with one matrix
cannot be told apart by a rank correlation. Aas, Czado, Frigessi and Bakken build multivariate dependence from bivariate blocks, a
route to families richer than the elliptical ones \citep{aas2009}.

These works give the families that fill the interior of the band. They do not say which one a bank should use for four risk types
of different nature, and for the pairs of risk types the matrix is a convention, as Section~\ref{sec:problem} states.

\subsection{Sharp bounds under dependence uncertainty}
\label{sec:rw-bounds}

When the marginals are known and the dependence is not, the value at risk of the sum lies between two sharp bounds. Makarov and
R\"uschendorf obtained the bounds on the distribution of a sum of two variables with fixed marginals \citep{makarov1981,
ruschendorf1982}. Embrechts, H\"oing and Puccetti studied the worst value at risk scenarios \citep{worstvar2005}, Puccetti and
R\"uschendorf proposed the rearrangement algorithm that computes the bounds for several risks \citep{puccetti2012}, and Embrechts,
Puccetti and R\"uschendorf applied it to the aggregation of risks and to the model uncertainty that the bounds measure
\citep{epr2013}. Hofert, Memartoluie, Saunders and Wirjanto improved the algorithm and examined its numerical challenges
\citep{hofert2015}, and Hofert describes the implementation of the algorithm as an example of computational risk management, from a basic version to a fast and reliable one \citep{hofert2020}. Wang, Peng and Yang
give bounds for monotone marginal densities \citep{wpy2013}, Puccetti, Wang and Wang relate the worst value at risk to the
expected shortfall \citep{pww2013}, Bernard, Jiang and Wang study the aggregation of risks under dependence uncertainty
\citep{bjw2014}, and Bernard, R\"uschendorf and Vanduffel narrow the bounds with a constraint on the variance of the sum
\citep{brv2015}. R\"uschendorf's monograph collects the theory \citep{ruschendorf2013}. Embrechts and co-authors use the bounds to
argue that small changes of a model can change the aggregate by a large amount \citep{academic2014}.

All of these deliver the two end points. They stop short of the use that the proposed protocol makes of them, which is to place the
end points on the same display as the copula figures that a bank actually reports and to test that the display is complete.

\subsection{Subadditivity, coherence and allocation}
\label{sec:rw-subadditivity}

Artzner, Delbaen, Eber and Heath fixed the axioms of a coherent risk measure, of which subadditivity is the one that expresses
diversification \citep{artzner1999}; Acerbi and Tasche showed that the expected shortfall is coherent \citep{acerbi2002}. The value
at risk is not, so the comonotonic sum is not an upper bound on it. Embrechts, Lambrigger and W\"uthrich give examples and
counter-examples of the aggregation of dependent heavy-tailed risks \citep{elw2008}, Danielsson and co-authors show that the value
at risk is subadditive in the tail for fat tails with a finite mean \citep{danielsson2013}, and Embrechts, Wang and Wang study how
robust a regulatory risk measure is to the aggregation \citep{ewang2015}. For the allocation of an aggregate to its components,
Tasche derives the Euler principle, under which the contribution of a component is the derivative of the aggregate risk with
respect to its weight \citep{tasche2008}.

The result used in what follows is the pair of facts that the value at risk of the sum can exceed the sum of the values at risk and
that the expected shortfall cannot, so that the expected shortfall of the stand-alone risks is an analytic ceiling for the value at
risk of the sum under any dependence (Section~\ref{sec:theory}).

\subsection{Model risk measures}
\label{sec:rw-modelrisk}

The uncertainty of the aggregate under alternative models has been given numbers. Cont proposed a framework for the model
uncertainty of the value of an instrument \citep{cont2006}; Kerkhof, Melenberg and Schumacher link model risk to capital reserves
\citep{kerkhof2010}; Barrieu and Scandolo define an additive and a relative measure of model risk over a set of models
\citep{barrieu2015}; Glasserman and Xu measure the worst case of a quantity over the alternatives within a budget of relative
entropy \citep{glasserman2014}; and Boucher, Danielsson, Kouontchou and Maillet estimate the model risk of risk forecasts
\citep{boucher2014}. These measures share one feature: each takes a set of alternative models as given and reports the spread of a
risk figure over it, so the number it returns is bounded by the set. A set built from one family of models gives a small number,
and a set that includes the sharp bounds gives the largest number that the marginals allow. The choice of the set is therefore the
substantive decision, and the band of the proposed protocol fixes it by rule: the set of models is the eight rows. The protocol
labels the measures information, as it labels every row but the headline.

\subsection{Where the proposed protocol stands}
\label{sec:rw-position}

Every ingredient of the protocol is published: the copula families, the sharp bounds, the subadditivity result, the Euler
allocation and the model-risk measures. The studies of Rosenberg and Schuermann and of Embrechts, Puccetti and R\"uschendorf
report the sensitivity of the aggregate to the dependence as a result of a study, and the supervisory documents report the
reluctance to rely on aggregation models. What this paper adds is a rule that composes the ingredients, one band on one set of
marginals at one horizon and one confidence level, with a single labelled headline, with the allocation and the model-risk
measures labelled alike, and with acceptance criteria that fail when the rule is broken; and a measurement of what the band
contains when the dependence structure is known because it was planted. The claim is limited to what the protocol does and what
the measurement shows: that no maintained toolkit or published practice ships the band beside every aggregate was not established
by a search beyond the sources cited here.
